\subsection{Centroid Decomposition}
\subsubsection{Centroid}
A centroid of a tree is defined as a node such that if removed, no connected component has  more than $\frac{N}{2}$ nodes.

We can find a centroid in a tree by selecting an arbitrary node in the tree. Then loop through its children's subtree. If the size of a child's subtree is more than $\frac{N}{2}$, traverse down that node. Otherwise, if no child's subtree has size more than $\frac{N}{2}$, then that node is the centroid.

\begin{lstlisting}
void fsz(int u, int p) {
    sz[u] = 1;
    for (auto v : adj[u]) if (v != p) fsz(v, u), sz[u] += sz[v];
}

int ct(int u, int p) {
    for (auto v : adj[u]) if (sz[v] * 2 > n) return ct(v, u);
    return u;
}
\end{lstlisting}

\noindent\textbf{Properties}
\begin{enumerate}
    \item If the centroid of a tree isn't unique, there are exactly two centroids which are adjacent. Removing the edge between them divides the tree into two connected components of the equal sizes.
    \item When a leaf is added to or removed from a tree, the centroid changes by at most one node.
    \item When two trees are connected, the centroid of the resulting tree lies on the path between the centroids of the original trees.
\end{enumerate}

\subsubsection{Centroid Decomposition}

Centroid Decomposition is a divide-and-conquer technique for trees, which recursively splits the tree by removing its centroid. The resulting decomposition tree has a height of at most $\mathcal{O}(\log N)$.

\noindent\textbf{Properties}
\begin{enumerate}
    \item The height of the decomposed tree is at most $\log{N}$
    \item Every path between $u$ and $v$ can be decomposed into two parts: $u$ to $\text{lca}_{ctd}(u, v)$ and $\text{lca}_{ctd}(u, v)$ to $v$. We can use this property for querying and updating. In the \textbf{Xania and Tree} problem, we update and query by traversing up to the $\log{N}$ parents for each node.
\end{enumerate}

\begin{lstlisting}
int fsz(int u, int p) {
    sz[u] = 1;
    for (auto v : adj[u]) if (v != p && !vis[v]) sz[u] += fsz(v, u);
    return sz[u];
}
 
int ct(int u, int p, int tsz) {
    for (auto v : adj[u]) if (v != p && !vis[v]) if (sz[v] * 2 > tsz) return ct(v, u, tsz);
    return u;
}
 
void ctd(int u, int p) {
    int C = ct(u, u, fsz(u, u));
    vis[C] = 1;
    if (p == -1) p = C;
    par[C] = p;
    for (auto v : adj[C]) if (!vis[v]) ctd(v, C);
}
\end{lstlisting}

\break
